\documentclass[../../main.tex]{subfiles}

\begin{document}

\section{闭集与极限点}\label{section:ch02-闭集与极限点}

我们已经掌握了少数几个拓扑空间的例子，以下我们着手引入一些与拓扑空间有关的基本概念。本节将讨论闭集、集合的闭包、极限点等概念，而这些又将导致我们去讨论拓扑空间中的一个公理，即 Hausdorff 公理。

\subsection*{闭集}

拓扑空间 $X$ 的一个子集 $A$，如果 $X - A$ 是开集，则称 $A$ 为一个闭集（closed set）。

\begin{example}\label{example:ch02-闭集1}
$\mathbb{R}$ 的子集 $[a, b]$ 是一个闭集，这是由于它的补
\begin{align*}
\mathbb{R} - [a, b] = (-\infty, a) \cup (b, +\infty)
\end{align*}
是一个开集。类似地，$[a, +\infty)$ 是一个闭集，这是由于它的补 $(-\infty, a)$ 是一个开集。这些事实正是我们使用“闭区间”、“闭射线”这类术语的原因。$\mathbb{R}$ 的子集 $[a, b)$ 既不是开集又不是闭集。
\end{example}

\begin{example}\label{example:ch02-闭集2}
在平面 $\mathbb{R}^{2}$ 上，集合
\begin{align*}
\{x \times y \mid x \geqslant 0 \quad \text{和} \quad y \geqslant 0\}
\end{align*}
是闭集，这是由于它的补是下面两个集合
\begin{align*}
(-\infty, 0) \times \mathbb{R} \quad \text{和} \quad \mathbb{R} \times (-\infty, 0)
\end{align*}
的并，而其中每一个都是 $\mathbb{R}$ 中两个开集的积，所以是 $\mathbb{R}^{2}$ 中的一个开集。
\end{example}

\begin{example}\label{example:ch02-闭集3}
对于集合 $X$ 的有限补拓扑，$X$ 自身及 $X$ 的所有有限集构成 $X$ 的闭集族。
\end{example}

\begin{example}\label{example:ch02-闭集4}
对于集合 $X$ 的离散拓扑，每一个集合都是开集，从而每一个集合也都是闭集。
\end{example}

\begin{example}\label{example:ch02-闭集5}
考虑实直线上具有子空间拓扑的子集
\begin{align*}
Y = [0, 1] \cup (2, 3),
\end{align*}
在这个空间中，由于 $[0, 1]$ 是 $\mathbb{R}$ 中的开集 $\left(-\frac{1}{2}, \frac{3}{2}\right)$ 与 $Y$ 的交，所以集合 $[0, 1]$ 是一个开集。类似地，$(2, 3)$ 作为 $Y$ 的子集也是一个开集，同时，作为 $\mathbb{R}$ 的子集也是开集。由于 $[0, 1]$ 及 $(2, 3)$ 在 $Y$ 中互为补，因此 $[0, 1]$ 和 $(2, 3)$ 作为 $Y$ 的子集都是闭集。
\end{example}

这些例子解开了数学家们心中的谜：“一个集合与一扇门究竟有什么不同呢？”现在可以说：“一扇门或是开，或是关，不可能是既开又关的。但是一个集合却可以是开的，可以是闭的，可以是既开又闭的，还可以是既不开又不闭的！”

空间 $X$ 的闭集族与开集族的性质十分类似。

\begin{theorem}[闭集的性质]\label{theorem:ch02-闭集性质}
设 $X$ 是一个拓扑空间。则下述结论成立：
\begin{enumerate}[(1)]
\item $\varnothing$ 和 $X$ 都是闭的.
\item 闭集的任意交都是闭的.
\item 闭集的有限并都是闭的.
\end{enumerate}
\end{theorem}

\begin{proof}
(1) 因为 $\varnothing$ 与 $X$ 的补分别是开集 $X$ 与 $\varnothing$，所以 $\varnothing$ 与 $X$ 都是闭集。

(2) 给定闭集的一个族 $\{A_{\alpha}\}_{\alpha \in J}$，应用 De Morgan 定律得到：
\begin{align*}
X - \bigcap_{\alpha \in J} A_{\alpha} = \bigcup_{\alpha \in J} (X - A_{\alpha}).
\end{align*}
根据定义，$X - A_{\alpha}$ 是开集，而等式右边是开集的任意并，所以它是开的，于是 $\bigcap A_{\alpha}$ 是闭的。

(3) 类似地，若对于 $i = 1, \cdots, n$，$A_{i}$ 为闭的，考虑以下等式：
\begin{align*}
X - \bigcup_{i=1}^{n} A_{i} = \bigcap_{i=1}^{n} (X - A_{i}).
\end{align*}
上式右边是开集的有限交，所以是开的，因此 $\bigcup A_{i}$ 是闭的。
\end{proof}

我们可以用满足本定理三条性质的集族（将其元素称之为“闭集”）替换开集族来刻画空间的拓扑。然后把开集定义成闭集的补并且完全像从前一样去做。由于这个做法不会给我们带来什么特别的方便，所以大多数数学家还是用开集来定义拓扑。

当涉及子空间时，对于“闭集”这个词的使用就要小心。设 $Y$ 是 $X$ 的一个子空间，我们称 $A$ 是 $Y$ 中的闭集，是指 $A$ 是 $Y$ 的一个子集，并且 $A$ 是 $Y$ 的子空间拓扑中的闭集（也就是说，$Y - A$ 是 $Y$ 中的开集）。有下面的定理。

\begin{theorem}[子空间闭集的刻画]\label{theorem:ch02-子空间闭集}
设 $Y$ 是 $X$ 的一个子空间。集合 $A$ 是 $Y$ 的一个闭集当且仅当 $A$ 是 $X$ 中的一个闭集与 $Y$ 的交。
\end{theorem}

\begin{proof}
假定 $A = C \cap Y$，其中 $C$ 为 $X$ 的一个闭集（参见\cref{figure:ch02-子空间闭集1}），那么 $X - C$ 是 $X$ 的一个开集。根据子空间拓扑的定义可见 $(X - C) \cap Y$ 是 $Y$ 的一个开集。而 $(X - C) \cap Y = Y - A$，因此 $Y - A$ 是 $Y$ 的一个开集，于是 $A$ 是 $Y$ 的一个闭集。反之，假定 $A$ 是 $Y$ 的一个闭集（参见\cref{figure:ch02-子空间闭集2}），那么 $Y - A$ 是 $Y$ 的一个开集。根据定义，$Y - A$ 等于 $X$ 的一个开集 $U$ 与 $Y$ 的交。$X - U$ 是 $X$ 的一个闭集并且 $A = Y \cap (X - U)$，因此 $A$ 等于 $X$ 的一个闭集与 $Y$ 的交。
\end{proof}

\par
\begin{figure}[htbp]
\centering
\includegraphics[width=0.35\textwidth]{ch02-fig17-1}
\caption{子空间中的闭集（一）}\label{figure:ch02-子空间闭集1}
\end{figure}

\par
\begin{figure}[htbp]
\centering
\includegraphics[width=0.35\textwidth]{ch02-fig17-2}
\caption{子空间中的闭集（二）}\label{figure:ch02-子空间闭集2}
\end{figure}

集合 $A$ 是空间 $X$ 的子空间 $Y$ 的一个闭集，那它可能是空间 $X$ 的闭集，也可能不是空间 $X$ 的闭集。如同开集的情形那样，集合 $A$ 是否是 $X$ 的闭集也有一个判别准则，以下定理留给读者自己证明。

\begin{theorem}[子空间闭集为全空间闭集的条件]\label{theorem:ch02-子空间闭集条件}
设 $Y$ 是 $X$ 的一个子空间。若 $A$ 是 $Y$ 的一个闭集并且 $Y$ 是 $X$ 的一个闭集，则 $A$ 是 $X$ 的一个闭集。
\end{theorem}

\subsection*{集合的闭包与内部}

拓扑空间 $X$ 的一个子集 $A$ 的内部（interior）定义为包含于 $A$ 的所有开集的并。$A$ 的闭包（closure）定义为包含着 $A$ 的所有闭集的交。$A$ 的内部记作 $\mathrm{Int}\,A$ 或 $\mathring{A}$，$A$ 的闭包记作 $\mathrm{Cl}\,A$ 或 $\overline{A}$。显然，$\mathring{A}$ 是开集，$\overline{A}$ 是闭集。因此有
\begin{align*}
\mathring{A} \subset A \subset \overline{A}.
\end{align*}
若 $A$ 是一个开集，则 $A = \mathring{A}$。若 $A$ 是一个闭集，则 $A = \overline{A}$。

集合的内部我们用得不多，但是集合的闭包却很重要。

如果涉及拓扑空间 $X$ 和一个子空间 $Y$，在取集合的闭包时需要十分小心。若 $A$ 是 $Y$ 的一个子集，那么 $A$ 在 $Y$ 中的闭包与 $A$ 在 $X$ 中的闭包一般并不相同。这时，我们仍然用记号 $\overline{A}$ 表示 $A$ 在 $X$ 中的闭包。而 $A$ 在 $Y$ 中的闭包可以藉助于 $\overline{A}$ 给出。见以下的定理。

\begin{theorem}[子空间中的闭包]\label{theorem:ch02-子空间闭包}
设 $Y$ 是 $X$ 的一个子空间，$A$ 是 $Y$ 的一个子集，$\overline{A}$ 表示 $A$ 在 $X$ 中的闭包。那么 $A$ 在 $Y$ 中的闭包等于 $\overline{A} \cap Y$。
\end{theorem}

\begin{proof}
用 $B$ 表示 $A$ 在 $Y$ 中的闭包。$\overline{A}$ 是 $X$ 中的闭集，根据\cref{theorem:ch02-子空间闭集}，$\overline{A} \cap Y$ 是 $Y$ 的闭集。因为 $\overline{A} \cap Y$ 包含 $A$，并且根据定义，$B$ 等于 $Y$ 中包含 $A$ 的所有闭子集的交，所以有 $B \subset (\overline{A} \cap Y)$。

另一方面，已知 $B$ 是 $Y$ 的一个闭集。根据\cref{theorem:ch02-子空间闭集}，存在 $X$ 中某闭集 $C$，使得 $B = C \cap Y$，则 $C$ 是 $X$ 中包含 $A$ 的一个闭集。又因为 $\overline{A}$ 是所有这种闭集的交，所以有 $\overline{A} \subset C$。于是 $(\overline{A} \cap Y) \subset (C \cap Y) = B$。
\end{proof}

集合闭包定义本身并没有向我们提供求出具体集合的闭包的简单方法，因为 $X$ 的闭集族就像开集族一样，往往是太多了。下面的定理给出了描述一个集合闭包的另一种方法，它仅仅用到 $X$ 的拓扑基，因而十分方便。

我们首先引进一个术语。所谓集合 $A$ 与集合 $B$ 相交（intersect），是指交 $A \cap B$ 不是空集。

\begin{theorem}[闭包的邻域判别]\label{theorem:ch02-闭包邻域}
设 $A$ 是拓扑空间 $X$ 的一个子集。
\begin{enumerate}[(a)]
\item $x \in \overline{A}$ 当且仅当每一个包含 $x$ 的开集 $U$ 与 $A$ 相交.
\item 假定 $X$ 的拓扑由一个基给出，则 $x \in \overline{A}$ 当且仅当含有 $x$ 的每一个基元素 $B$ 与 $A$ 相交.
\end{enumerate}
\end{theorem}

\begin{proof}
考虑论断 (a)，它是 $P \Leftrightarrow Q$ 的形式。我们对于每一个蕴涵关系都考虑它的逆否命题，从而得到一个等价的论断（非 $P$）$\Leftrightarrow$（非 $Q$）。正式写出来是：
\begin{align*}
x \notin \overline{A} \Longleftrightarrow \text{存在一个包含 } x \text{ 的开集 } U \text{，它与 } A \text{ 无交}.
\end{align*}
把定理写成这种形式就容易证明了。若 $x \notin \overline{A}$，则包含 $x$ 的一个开集 $U = X - \overline{A}$ 与 $A$ 无交。反之，若存在包含 $x$ 的一个开集 $U$ 与 $A$ 无交，则 $X - U$ 是包含 $A$ 的一个闭集。根据闭包 $\overline{A}$ 的定义，$X - U$ 必定包含 $\overline{A}$。因此 $x$ 不可能在 $\overline{A}$ 中。

论断 (b) 的证明是容易的。如果每一个包含 $x$ 的开集都与 $A$ 相交，那么每一个包含 $x$ 的基元素 $B$ 也与 $A$ 相交，这是因为 $B$ 是开集。反之，如果每一个包含 $x$ 的基元素都与 $A$ 相交，那么每一个包含 $x$ 的开集 $U$ 也与 $A$ 相交，这是因为 $U$ 包含一个含有 $x$ 的基元素。
\end{proof}

数学家们往往在这里使用一个专门术语，他们把“$U$ 是包含 $x$ 的一个开集”这句话简单地说成——$U$ 是 $x$ 的一个邻域（neighborhood）。

使用这个术语，上面定理的前半段就可以写成：

若 $A$ 是拓扑空间 $X$ 的一个子集，则 $x \in \overline{A}$ 当且仅当 $x$ 的每一个邻域与 $A$ 相交。

\begin{example}\label{example:ch02-闭包例1}
设 $X$ 为实直线 $\mathbb{R}$。若 $A = (0, 1]$，则 $\overline{A} = [0, 1]$。事实上，0 的每一个邻域都与 $A$ 相交，而 $[0, 1]$ 以外的任意点都有不与 $A$ 相交的邻域。对于 $X$ 的下列子集，我们可以类似地证明：

若 $B = \left\{\frac{1}{n} \mid n \in \mathbb{Z}_{+}\right\}$，则 $\overline{B} = \{0\} \cup B$。若 $C = \{0\} \cup (1, 2)$，则 $\overline{C} = \{0\} \cup [1, 2]$。若 $\mathbb{Q}$ 为有理数集，$\overline{\mathbb{Q}} = \mathbb{R}$。若 $\mathbb{Z}_{+}$ 为正整数集，则 $\overline{\mathbb{Z}}_{+} = \mathbb{Z}_{+}$。若 $\mathbb{R}_{+}$ 为正实数集，则 $\mathbb{R}_{+}$ 的闭包是 $\mathbb{R}_{+} \cup \{0\}$。（这就是我们在第\cref{section:ch02-子空间拓扑}节中采用记号 $\overline{\mathbb{R}}_{+}$ 表示 $\mathbb{R}_{+} \cup \{0\}$ 的理由。）
\end{example}

\begin{example}\label{example:ch02-闭包例2}
考虑实直线 $\mathbb{R}$ 的子空间 $Y = (0, 1]$。$A = \left(0, \frac{1}{2}\right)$ 是 $Y$ 的一个子集，它在 $\mathbb{R}$ 中的闭包是 $\left[0, \frac{1}{2}\right]$，它在 $Y$ 中的闭包是 $\left[0, \frac{1}{2}\right] \cap Y = \left(0, \frac{1}{2}\right]$。
\end{example}

有些数学家在不同意义上使用“邻域”一词。他们称 $A$ 为 $x$ 的一个邻域，仅仅是指 $A$ 包含一个含有 $x$ 的开集。我们不用这种说法。

\subsection*{极限点}

还有一种描述集合闭包的方法，它要用到我们下面将要讨论的极限点这样一个重要的概念。

若 $A$ 是拓扑空间 $X$ 的一个子集，$x$ 是 $X$ 的一个点，如果 $x$ 的任意一个邻域与 $A$ 的交含有异于 $x$ 的点，我们称 $x$ 为 $A$ 的一个极限点（limit point）（或“聚点”（cluster point, point of accumulation））。换言之，如果 $x$ 属于 $A - \{x\}$ 的闭包，则 $x$ 便是 $A$ 的一个极限点。这个点 $x$ 可以在 $A$ 中，也可以不在 $A$ 中，在此定义中 $x$ 是否在 $A$ 中这一点无关紧要。

\begin{example}\label{example:ch02-极限点例}
考虑实直线 $\mathbb{R}$。若 $A = (0, 1]$，则 0 是 $A$ 的一个极限点，点 $\frac{1}{2}$ 也是 $A$ 的一个极限点。事实上，区间 $[0, 1]$ 中的任何一个点都是 $A$ 的极限点，而除此之外的 $\mathbb{R}$ 的点都不再是 $A$ 的极限点。

若 $B = \left\{\frac{1}{n} \mid n \in \mathbb{Z}_{+}\right\}$，则 0 是 $B$ 的唯一极限点。$\mathbb{R}$ 中除 0 之外的任何点都有一个邻域，或与 $B$ 无交，或与 $B$ 只交于一点 $x$。若 $C = \{0\} \cup (1, 2)$，则区间 $[1, 2]$ 中的点都是 $C$ 的极限点。若 $\mathbb{Q}$ 是有理数集，则 $\mathbb{R}$ 的任何点都是 $\mathbb{Q}$ 的极限点。若 $\mathbb{Z}_{+}$ 是正整数集，则 $\mathbb{R}$ 的任何点都不是 $\mathbb{Z}_{+}$ 的极限点。若 $\mathbb{R}_{+}$ 是正实数集，则 $\{0\} \cup \mathbb{R}_{+}$ 中的每一个点都是 $\mathbb{R}_{+}$ 的极限点。
\end{example}

参照\cref{example:ch02-闭包例1}与\cref{example:ch02-极限点例}，可以发现一个集合的闭包与它的极限点之间的关系。这就是下面的定理。

\begin{theorem}[闭包与极限点]\label{theorem:ch02-闭包极限点}
设 $A$ 是拓扑空间 $X$ 的一个子集，$A'$ 是 $A$ 的所有极限点的集合，则
\begin{align*}
\overline{A} = A \cup A'.
\end{align*}
\end{theorem}

\begin{proof}
若 $x \in A'$，那么 $x$ 的每一个邻域与 $A$ 的交有异于 $x$ 的点，这样根据\cref{theorem:ch02-闭包邻域}可见 $x$ 属于 $\overline{A}$，因此 $A' \subset \overline{A}$。然而根据定义有 $A \subset \overline{A}$，所以 $A \cup A' \subset \overline{A}$。

为了证明反过来的包含关系成立，令 $x$ 为 $\overline{A}$ 的一个点并证明 $x \in A \cup A'$。如果 $x$ 是 $A$ 的一个点，那么显然 $x \in A \cup A'$。现在设 $x$ 不属于 $A$，因为 $x \in \overline{A}$，那么 $x$ 的每一个邻域 $U$ 都与 $A$ 相交，由于 $x \notin A$，集合 $U$ 必与 $A$ 交于异于 $x$ 的一个点，因此 $x \in A'$。于是 $x \in A \cup A'$。
\end{proof}

\begin{corollary}[闭集包含其极限点]\label{corollary:ch02-闭集极限点}
拓扑空间的一个子集是闭集当且仅当它包含其所有极限点。
\end{corollary}

\begin{proof}
集合 $A$ 是闭集当且仅当 $A = \overline{A}$，而 $A = \overline{A}$ 当且仅当 $A' \subset A$。
\end{proof}

\subsection*{Hausdorff 空间}

在我们研究更为一般的拓扑空间的时候，如果仅凭有关实直线和平面上开集、闭集、极限点的经验，就会出现错误。在 $\mathbb{R}$ 和 $\mathbb{R}^{2}$ 中，每一个单点集 $\{x_{0}\}$ 都是闭集。这个事实是容易证明的，事实上，每一个异于 $x_{0}$ 的点都有一个邻域不与 $\{x_{0}\}$ 相交。因此 $\{x_{0}\}$ 的闭包就是 $\{x_{0}\}$。但是这一点对于任意拓扑空间却不成立。例如，我们在开始时就讲过的一个拓扑空间，如\cref{figure:ch02-非Hausdorff}所示的三点集 $\{a, b, c\}$ 的拓扑。对于这个空间，单点集 $\{b\}$ 就不是闭集，这是由于它的补不是开集的缘故。

\par
\begin{figure}[htbp]
\centering
\includegraphics[width=0.3\textwidth]{ch02-fig17-3}
\caption{非 Hausdorff 空间}\label{figure:ch02-非Hausdorff}
\end{figure}

类似地，在我们处理更为一般的拓扑空间时，如果仅凭有关 $\mathbb{R}$ 和 $\mathbb{R}^{2}$ 中收敛序列的经验，也可能会导致错误。对于任意拓扑空间 $X$，设 $x_{1}, x_{2}, \cdots$ 是 $X$ 中的一个序列并且 $x \in X$。如果对于 $x$ 的任何一个邻域 $U$，存在一个正整数 $N$，使得当 $n > N$ 时，$x_{n} \in U$，则称 $X$ 中的序列 $x_{1}, x_{2}, \cdots$ 收敛（converge）到点 $x$。在 $\mathbb{R}$ 和 $\mathbb{R}^{2}$ 中，一个序列最多收敛到一个点，但在一般拓扑空间中，一个序列可能收敛到多个点。比如，在\cref{figure:ch02-非Hausdorff}所示的拓扑空间中，序列 $x_{n} = b$ 不仅收敛到点 $b$，它也收敛到点 $a$ 和点 $c$！

使得单点集不是闭集的拓扑，以及使得一个序列可能收敛到多个点的拓扑，被许多数学家认为是非正常的拓扑。这些并不能引起数学家们多大的兴趣，因为它们在数学的其他分支中很少出现。如果这种例子被容许的话，那对于拓扑空间能够证明的定理就是很有限的。因此常常需要一些附加条件，使得上述例子相应的情况不会出现，从而使我们所讨论的空间更加接近几何的直观。这个条件是数学家 Felix Hausdorff 给出的，因此数学家们就以他的名字来命名。

\begin{definition}[Hausdorff 空间]\label{definition:ch02-Hausdorff}
如果对于拓扑空间 $X$ 中任意两个不同的点 $x_{1}$ 和 $x_{2}$，分别存在 $x_{1}$ 和 $x_{2}$ 的邻域 $U_{1}$ 和 $U_{2}$ 使得这两个邻域无交，则称 $X$ 为一个 Hausdorff 空间（Hausdorff space）。
\end{definition}

\begin{theorem}[Hausdorff 空间中有限集是闭集]\label{theorem:ch02-Hausdorff有限闭}
Hausdorff 空间中的任何有限集都是闭的。
\end{theorem}

\begin{proof}
我们只要证明任何一个单点集 $\{x_{0}\}$ 是闭集就行了。设 $x$ 为 $X$ 中异于 $x_{0}$ 的一个点，那么 $x$ 和 $x_{0}$ 分别有无交的邻域 $U$ 和 $V$。于是 $U$ 与 $\{x_{0}\}$ 无交，$x$ 就不属于 $\{x_{0}\}$ 的闭包，因此 $\{x_{0}\}$ 的闭包就是 $\{x_{0}\}$，所以 $\{x_{0}\}$ 是闭集。
\end{proof}

有限点集是闭集这一条件实际上是比 Hausdorff 条件更弱的条件。例如，实直线 $\mathbb{R}$ 关于有限补拓扑并不是一个 Hausdorff 空间，但在此空间中有限点集是闭集。有限点集是闭集的条件也有一个名字，叫做 $T_{1}$ 公理（$T_{1}$ axiom）（我们将在第 4 章对其做进一步说明）。本书仅附有少量的有关 $T_{1}$ 公理的练习及以下的一个相关定理。

\begin{theorem}[极限点的邻域判别（$T_{1}$ 公理）]\label{theorem:ch02-T1极限点}
设拓扑空间 $X$ 满足 $T_{1}$ 公理，$A$ 是 $X$ 的一个子集。则点 $x$ 是 $A$ 的极限点当且仅当 $x$ 的每一个邻域与 $A$ 的交是一个无限集。
\end{theorem}

\begin{proof}
设 $x$ 的任意邻域与 $A$ 相交于无穷多个点，则必交于异于 $x$ 的一个点，于是 $x$ 是 $A$ 的一个极限点。

反之，设 $x$ 是 $A$ 的一个极限点，$x$ 的某邻域 $U$ 与 $A$ 仅相交于有限多个点，于是 $U$ 与 $A - \{x\}$ 也只相交于有限多个点。设 $\{x_{1}, \cdots, x_{m}\}$ 等于 $U \cap (A - \{x\})$，因为有限集 $\{x_{1}, \cdots, x_{m}\}$ 是闭集，所以 $X - \{x_{1}, \cdots, x_{m}\}$ 是 $X$ 的一个开集。因此
\begin{align*}
U \cap (X - \{x_{1}, \cdots, x_{m}\})
\end{align*}
就是与 $A - \{x\}$ 无交的 $x$ 的一个邻域，这与 $x$ 是 $A$ 的极限点的假定相矛盾。
\end{proof}

我们对 $T_{1}$ 公理不是很感兴趣，其原因在于有关拓扑的许多有趣的定理不仅需要 $T_{1}$ 公理，而且需要强于它的 Hausdorff 公理。此外，那些对于数学家而言重要的空间大多都是 Hausdorff 空间。下面的两个定理是以上说法的佐证。

\begin{theorem}[Hausdorff 空间中序列极限唯一]\label{theorem:ch02-Hausdorff极限唯一}
若 $X$ 是一个 Hausdorff 空间，则 $X$ 中的一个序列最多收敛到一个点。
\end{theorem}

\begin{proof}
假设 $X$ 中点的序列 $x_{n}$ 收敛到 $x$。对于 $y \neq x$，记 $U$ 和 $V$ 分别是 $x$ 和 $y$ 的两个无交邻域，由于 $U$ 包含着除有限个之外所有的 $x_{n}$，所以 $V$ 一定不再会这样。因此，$x_{n}$ 不会收敛到 $y$。
\end{proof}

Hausdorff 空间 $X$ 中的序列 $x_{n}$ 收敛到 $X$ 中的点 $x$，我们通常记为 $x_{n} \rightarrow x$，并称 $x$ 为序列 $x_{n}$ 的极限（limit）。

下述定理的证明留给读者。

\begin{theorem}[Hausdorff 空间的封闭性质]\label{theorem:ch02-Hausdorff封闭}
每一个具有序拓扑的全序集是一个 Hausdorff 空间。两个 Hausdorff 空间的积是一个 Hausdorff 空间。Hausdorff 空间的子空间是一个 Hausdorff 空间。
\end{theorem}

Hausdorff 条件一般被认为是加到拓扑空间上的一个很合适的附加条件。事实上，在拓扑学基本教程中，有些数学家甚至在一开始就加上这个条件，而对不满足 Hausdorff 条件的空间不加研究。我们不准备这么办，但是在证明需要的时候，总是假定有 Hausdorff 条件，除非它对于结果的应用范围产生了令人不能满意的苛刻限制。

Hausdorff 条件是能够加到拓扑空间上的许多条件中的一个。人们每加上这样一个条件，就可以证明更强的定理，但也限制了适用这些定理的空间类。从拓扑学创立以来，相当一部分研究工作就是在寻求这样一些条件，这些条件比较强，但能使人们在满足这些条件的空间上证明一些有意义的定理；然而这些条件又不能太强，以免导致其结论的应用受到苛刻的限制。

我们将在后面两章研究这样一些条件，已知的 Hausdorff 条件以及 $T_{1}$ 公理只不过是我们熟悉的通称为分离性公理的一个系列条件中的两个。其他的条件包括可数性公理、紧致性以及连通性等，你将会发现，它们中的一些确实是相当强的限制。

\subsection*{习题}

\begin{exercise}\label{exercise:ch02-闭集与极限点-1}
设 $\mathcal{C}$ 是集合 $X$ 的子集的一个族。假定 $\varnothing$ 与 $X$ 在 $\mathcal{C}$ 中，并且 $\mathcal{C}$ 中元素的有限并及任意交也在 $\mathcal{C}$ 中。证明集族 $\mathcal{T} = \{X - C \mid C \in \mathcal{C}\}$ 是 $X$ 的一个拓扑。
\end{exercise}

\begin{exercise}\label{exercise:ch02-闭集与极限点-2}
证明：若 $A$ 是 $Y$ 的闭集并且 $Y$ 是 $X$ 的闭集，则 $A$ 是 $X$ 的闭集。
\end{exercise}

\begin{exercise}\label{exercise:ch02-闭集与极限点-3}
证明：若 $A$ 是 $X$ 的闭集并且 $B$ 是 $Y$ 的闭集，则 $A \times B$ 是 $X \times Y$ 的闭集。
\end{exercise}

\begin{exercise}\label{exercise:ch02-闭集与极限点-4}
证明：若 $U$ 是 $X$ 的开集并且 $A$ 是 $X$ 的闭集，则 $U - A$ 是 $X$ 的开集，并且 $A - U$ 是 $X$ 的闭集。
\end{exercise}

\begin{exercise}\label{exercise:ch02-闭集与极限点-5}
设 $X$ 是具有序拓扑的全序集。证明：$\overline{(a, b)} \subset [a, b]$。等号在什么条件下成立？
\end{exercise}

\begin{exercise}\label{exercise:ch02-闭集与极限点-6}
设 $A$，$B$ 和 $A_{\alpha}$ 都是空间 $X$ 的子集。证明：
\begin{enumerate}[(a)]
\item 若 $A \subset B$，则 $\overline{A} \subset \overline{B}$.
\item $\overline{A \cup B} = \overline{A} \cup \overline{B}$.
\item $\overline{\bigcup A_{\alpha}} \supset \bigcup \overline{A_{\alpha}}$。举出等号不成立的例子.
\end{enumerate}
\end{exercise}

\begin{exercise}\label{exercise:ch02-闭集与极限点-7}
评述以下对于 $\overline{\bigcup A_{\alpha}} \subset \bigcup \overline{A_{\alpha}}$ 的证明：若 $\{A_{\alpha}\}$ 是 $X$ 中的集合的一个族，$x \in \overline{\bigcup A_{\alpha}}$，则 $x$ 的任意邻域 $U$ 与 $\bigcup A_{\alpha}$ 相交，因此 $U$ 必与某一个 $A_{\alpha}$ 相交，所以 $x$ 必定属于某一个 $A_{\alpha}$ 的闭包。从而 $x \in \bigcup \overline{A_{\alpha}}$。
\end{exercise}

\begin{exercise}\label{exercise:ch02-闭集与极限点-8}
设 $A$，$B$ 和 $A_{\alpha}$ 都是空间 $X$ 的子集。判断下面哪些等号成立。如果等号不成立，判断哪一种包含关系（$\supset$ 或 $\subset$）成立。
\begin{enumerate}[(a)]
\item $\overline{A \cap B} = \overline{A} \cap \overline{B}$.
\item $\overline{\bigcap A_{\alpha}} = \bigcap \overline{A_{\alpha}}$.
\item $\overline{A - B} = \overline{A} - \overline{B}$.
\end{enumerate}
\end{exercise}

\begin{exercise}\label{exercise:ch02-闭集与极限点-9}
设 $A \subset X$ 并且 $B \subset Y$。证明：在空间 $X \times Y$ 中 $\overline{A \times B} = \overline{A} \times \overline{B}$。
\end{exercise}

\begin{exercise}\label{exercise:ch02-闭集与极限点-10}
证明：任何序拓扑都是 Hausdorff 的。
\end{exercise}

\begin{exercise}\label{exercise:ch02-闭集与极限点-11}
证明：两个 Hausdorff 空间的积是 Hausdorff 的。
\end{exercise}

\begin{exercise}\label{exercise:ch02-闭集与极限点-12}
证明：Hausdorff 空间的子空间是 Hausdorff 的。
\end{exercise}

\begin{exercise}\label{exercise:ch02-闭集与极限点-13}
证明：$X$ 是一个 Hausdorff 空间当且仅当其对角线（diagonal）$\Delta = \{x \times x \mid x \in X\}$ 是 $X \times X$ 中的一个闭集。
\end{exercise}

\begin{exercise}\label{exercise:ch02-闭集与极限点-14}
在 $\mathbb{R}$ 的有限补拓扑空间中，序列 $x_{n} = \frac{1}{n}$ 收敛到哪一点或哪些点？
\end{exercise}

\begin{exercise}\label{exercise:ch02-闭集与极限点-15}
证明 $T_{1}$ 公理等价于：$X$ 中的每两个不同的点各自有一个邻域不包含另一点。
\end{exercise}

\begin{exercise}\label{exercise:ch02-闭集与极限点-16}
考虑在第\cref{section:ch02-拓扑的基}节习题 7 中所给的 $\mathbb{R}$ 的 5 个拓扑。

(a) 对于每一个拓扑，确定集合 $K = \{1/n \mid n \in \mathbb{Z}_{+}\}$ 的闭包。

(b) 其中的哪些拓扑满足 Hausdorff 公理，哪些满足 $T_{1}$ 公理？
\end{exercise}

\begin{exercise}\label{exercise:ch02-闭集与极限点-17}
考虑 $\mathbb{R}$ 的下限拓扑及由第\cref{section:ch02-拓扑的基}节习题 8 中的基 $\mathcal{C}$ 所给出的拓扑。对于上述两种拓扑确定区间 $A = (0, \sqrt{2})$ 及 $B = (\sqrt{2}, 3)$ 的闭包。
\end{exercise}

\begin{exercise}\label{exercise:ch02-闭集与极限点-18}
确定有序矩形中下列子集的闭包：
\begin{align*}
A = \{(1/n) \times 0 \mid n \in \mathbb{Z}_{+}\},
\end{align*}
\begin{align*}
B = \{(1 - 1/n) \times \frac{1}{2} \mid n \in \mathbb{Z}_{+}\},
\end{align*}
\begin{align*}
C = \{x \times 0 \mid 0 < x < 1\},
\end{align*}
\begin{align*}
D = \left\{x \times \frac{1}{2} \mid 0 < x < 1\right\},
\end{align*}
\begin{align*}
E = \left\{\frac{1}{2} \times y \mid 0 < y < 1\right\}.
\end{align*}
\end{exercise}

\begin{exercise}\label{exercise:ch02-闭集与极限点-19}
设 $A \subset X$，$A$ 的边界（boundary）定义为
\begin{align*}
\mathrm{Bd}\,A = \overline{A} \cap (\overline{X - A}).
\end{align*}
(a) 证明：$\mathrm{Int}\,A$ 与 $\mathrm{Bd}\,A$ 无交，并且 $\overline{A} = \mathrm{Int}\,A \cup \mathrm{Bd}\,A$。

(b) 证明：$\mathrm{Bd}\,A = \varnothing \Longleftrightarrow A$ 既开又闭。

(c) 证明：$U$ 是开集 $\Longleftrightarrow \mathrm{Bd}\,U = \overline{U} - U$。

(d) 若 $U$ 是开集，那么 $U = \mathrm{Int}(\overline{U})$ 成立吗？验证你的结论。
\end{exercise}

\begin{exercise}\label{exercise:ch02-闭集与极限点-20}
求出 $\mathbb{R}^{2}$ 的下列子集的每一个的边界及内部：
\begin{enumerate}[(a)]
\item $A = \{x \times y \mid y = 0\}$.
\item $B = \{x \times y \mid x > 0 \quad \text{并且} \quad y \neq 0\}$.
\item $C = A \cup B$.
\item $D = \{x \times y \mid x \text{ 是有理数}\}$.
\item $E = \{x \times y \mid 0 < x^{2} - y^{2} \leqslant 1\}$.
\item $F = \{x \times y \mid x \neq 0 \quad \text{并且} \quad y \leqslant 1/x\}$.
\end{enumerate}
\end{exercise}

\begin{exercise}\label{exercise:ch02-闭集与极限点-21}
(Kuratowski) 考虑拓扑空间 $X$ 的所有子集 $A$ 的族。闭包运算 $A \rightarrow \overline{A}$ 与取补运算 $A \rightarrow X - A$ 是从这个族到它自身的两个函数。

(a) 证明：逐次进行这两种运算，对于给定的集合 $A$，最多可以组成 14 个不同的集合。

(b) 对于具有通常拓扑的 $\mathbb{R}$，举出它的一个子集 $A$，由它恰好得到 14 个不同的集合。
\end{exercise}

\end{document}